\chapter{Umsetzung}
\label{kap:umsetzung}

Das in Kapitel~\ref{kap:modell} entwickelte Modell wird auf zwei Wegen
umgesetzt: als \Index{Modellierungssprache} in MathProg, die ein
Standardlöser bearbeitet, und als eigenes Verfahren in Python. Beide
Umsetzungen arbeiten auf derselben Instanz. Das ist keine
Selbstverständlichkeit und wird in Abschnitt~\ref{sec:einequelle}
begründet.

\section{Eine Quelle für die Daten}
\label{sec:einequelle}

Wer zwei Verfahren vergleicht, vergleicht in Wahrheit oft zwei
Datensätze. Werden die Zahlen für den Löser und für das eigene Programm
getrennt gepflegt, laufen sie früher oder später auseinander, und der
Vergleich misst dann den Unterschied der Eingaben statt den der
Verfahren.

Deshalb liegt die Instanz genau einmal vor, nämlich in
\texttt{Code/daten.py}. Aus dieser Quelle wird die Datendatei für
\ac{glpk} erzeugt; dasselbe Modul liest das Python-Verfahren direkt ein.
Die Bearbeitungszeit ergibt sich in beiden Fällen aus derselben Formel

\begin{equation}
  p_{ij} = \frac{g_j}{e_i(q_j)},
  \label{eq:dauer}
\end{equation}

wobei $g_j$ die \Index{Grundzeit} des Arbeitsgangs $j$ ist, $q_j$ die
benötigte Qualifikation und $e_i(q_j)$ die Effizienz des Mitarbeiters
$i$ in dieser Qualifikation. Ist $e_i(q_j)$ nicht definiert, fehlt die
Qualifikation, und das Paar $(i,j)$ existiert nicht.

Dieselbe Quelle speist auch die Abbildungen, die Tabellen und die
Zahlen im Fließtext dieser Arbeit. Alle Zahlenwerte in diesem Text sind
Makros, die aus den Rechenergebnissen erzeugt werden. Eine Zahl, die
hier steht, kann deshalb nicht von der Rechnung abweichen, auf die sie
sich beruft.

\section{Das Modell in MathProg}
\label{sec:mathprog}

\Index{MathProg} (auch GMPL) ist die Modellierungssprache von
\ac{glpk} \parencite{glpk}. Sie beschreibt das Modell deklarativ: Mengen,
Parameter, Variablen, Zielfunktion, Nebenbedingungen. Wie das Problem
gelöst wird, steht nicht im Modell -- das ist Sache des Lösers.

Die Mengen und die Zulässigkeit werden so erklärt:

\begin{lstlisting}[language=mathprog,caption={Mengen, Parameter und Variablen},label={lst:modmengen}]
set MITARBEITER;
set AUFGABEN;
set ZULAESSIG within {MITARBEITER, AUFGABEN};

param p{(i,j) in ZULAESSIG} > 0;

var x{(i,j) in ZULAESSIG}, binary;
var Cmax >= 0;
\end{lstlisting}

Entscheidend ist die dritte Zeile. \texttt{ZULAESSIG} ist eine Teilmenge
des Kreuzprodukts, und alle folgenden Größen sind nur über dieser
Teilmenge erklärt. Eine Variable $x_{ij}$ für ein Paar ohne passende
Qualifikation entsteht damit gar nicht erst. Die Alternative -- alle
Paare anlegen und die unzulässigen mit einer großen Zahl bestrafen --
wäre nicht nur größer, sondern brächte auch die in
Abschnitt~\ref{sec:eigenschaften} beschriebene Schwächung der
Relaxation mit sich.

Zielfunktion und Nebenbedingungen folgen unmittelbar der Formulierung
aus Kapitel~\ref{kap:modell}:

\begin{lstlisting}[language=mathprog,caption={Zielfunktion und Nebenbedingungen},label={lst:modnb}]
minimize Gesamtdauer: Cmax;

s.t. Vergabe{j in AUFGABEN}:
  sum{(i,j) in ZULAESSIG} x[i,j] = 1;

s.t. Auslastung{i in MITARBEITER}:
  sum{(i,j) in ZULAESSIG} p[i,j] * x[i,j] <= Cmax;
\end{lstlisting}

Das vollständige Modell steht in Anhang~\ref{anh:mathprog}. Aufgerufen
wird es mit

\begin{lstlisting}[language=bash,numbers=none]
glpsol --math erpolino.mod --data erpolino.dat
\end{lstlisting}

\ac{glpk} löst die Instanz in \MilpSekunden{} Sekunden und meldet
\texttt{INTEGER OPTIMAL SOLUTION FOUND}. Diese Meldung ist mehr als ein
Ergebnis: Sie ist ein Beweis. Der Löser hat gezeigt, dass keine bessere
Zuordnung existiert, nicht nur, dass er keine gefunden hat.

\section{Das Verfahren in Python}
\label{sec:python}

Das Verfahren folgt dem klassischen Ablauf
\parencite{kirkpatrick1983optimization}: Von einer Startlösung ausgehend
wird wiederholt eine Nachbarlösung erzeugt und angenommen oder
verworfen. Als Startlösung dient die \ac{lpt}-Regel.

\subsection{Nachbarschaft}

Zwei Operatoren erzeugen Nachbarlösungen, wie in
Abbildung~\ref{abb:nachbarschaft} dargestellt. Das \Index{Verschieben}
nimmt dem am stärksten ausgelasteten Mitarbeiter eine Aufgabe ab und
gibt sie einem anderen qualifizierten Mitarbeiter. Das \Index{Tauschen}
kreuzt zwei Aufgaben zwischen ihren Mitarbeitern, sofern beide dafür
qualifiziert sind.

Der zweite Operator ist nicht entbehrlich. Sind zwei Mitarbeiter beide
ausgelastet, führt jedes Verschieben zwangsläufig zu einer
Verschlechterung an anderer Stelle; nur ein Tausch kann die Last
umverteilen, ohne sie zu erhöhen. Ohne den Tausch bleibt das Verfahren
an solchen Konstellationen hängen.

\subsection{Die Energie und das Plateauproblem}
\label{sec:plateau}

Naheliegend wäre, die Zielgröße selbst als \Index{Energie} zu verwenden.
Das funktioniert schlecht, und der Grund ist lehrreich. Der
\Index{Makespan} ist ein Maximum: Er ändert sich nur, wenn gerade der
Engpass betroffen ist. Jeder Zug, der die übrigen Mitarbeiter betrifft,
ist energieneutral. Das Verfahren läuft über weite \Index{Plateaus} und
erhält kein Signal, welche Richtung sich lohnt.

Abhilfe schafft ein zusätzlicher Term, der schon dann sinkt, wenn die
Last gleichmäßiger verteilt wird:

\begin{equation}
  E(x) = \max_{i \in I} a_i(x)
       + \lambda \sqrt{\sum_{i \in I} a_i(x)^2},
  \qquad
  a_i(x) = \sum_{j:(i,j) \in Z} p_{ij}\, x_{ij}.
  \label{eq:energie}
\end{equation}

Der zweite Summand ist die euklidische Norm des
\Index{Auslastungsvektors}. Er macht die Landschaft begehbar, ohne die
Zielgröße zu ersetzen: Das Gewicht $\lambda = 0{,}002$ ist so klein
gewählt, dass der Term die Rangfolge zweier Lösungen mit
unterschiedlichem Makespan nicht umdrehen kann. Die Auslastungen liegen
bei rund 200~Minuten, der Term damit bei etwa 1,2; ein Unterschied im
Makespan von einer Minute wiegt schwerer.

Bewertet wird während des Laufs über $E$, berichtet und gespeichert wird
stets der echte Makespan. Was diese Änderung bewirkt, zeigt
Abschnitt~\ref{sec:energiewirkung}: Der Mittelwert über 20 Läufe sinkt
von \SaRohMittel{} auf \SaMittel{} Minuten, und erst mit dem
Glättungsterm wird das Optimum überhaupt erreicht.

\subsection{Abkühlung}

Die Temperatur fällt geometrisch, $T_{k+1} = \alpha\,T_k$, von
$T_{\text{start}} = 25$ auf $T_{\text{ende}} = 0{,}05$ über
\num{60000} Schritte. Der Faktor ergibt sich daraus zu

\begin{equation}
  \alpha = \left(\frac{T_{\text{ende}}}{T_{\text{start}}}\right)^{1/K}
  \approx 0{,}99990.
\end{equation}

Abbildung~\ref{abb:abkuehlung} zeigt den Verlauf und die daraus
folgende Annahmewahrscheinlichkeit für eine Verschlechterung. Zu Beginn
werden auch deutliche Verschlechterungen angenommen, gegen Ende kaum
noch: Das Verfahren geht von der Erkundung in die Feinsuche über.

\begin{figure}[tb]
  \centering
  % Geometrische Abkühlung und die daraus folgende Annahmewahrscheinlichkeit.
%
% Echtes MetaPost, aber über luamplib direkt aus LuaLaTeX heraus gerechnet:
% kein eigener mpost-Lauf nötig, und die Beschriftungen erben Schrift und
% Mathematiksatz des Dokuments. Die Farben kommen aus nxtfarben.mp, das
% Code/farben.py aus nxtlatexcolors.sty erzeugt -- dieselbe Quelle wie für
% LaTeX, TikZ und matplotlib.
%
% Zwei Eigenheiten von MetaPost, über die man hier stolpert:
%  1. Ein Name darf nicht auf einer Ziffer enden -- »T0« wäre »T« mit dem
%     Suffix 0. Deshalb Tstart.
%  2. Zahlen sind auf rund 32768 begrenzt, alpha^60000 läuft über. Der
%     Verlauf wird daher über die normierte Laufvariable t in [0,1]
%     gerechnet: T/Tstart = exp(-c*t) mit c = ln(Tstart/Tende).
\begin{mplibcode}[alt={Kurvendiagramm über die Iterationen eines Laufs:
  Die Temperatur fällt geometrisch von der Starttemperatur gegen null. Die
  Wahrscheinlichkeit, eine Verschlechterung anzunehmen, bleibt zunächst
  hoch und fällt dann ebenfalls auf null.}]
% Pfad relativ zum Hauptdokument, nicht zu dieser Datei.
input Abbildungen/nxtfarben;

numeric b, h;
b := 250;
h := 105;

% Abkühlung im Verfahren: Tstart = 25, Tende = 0,05, also c = ln(500).
numeric c;
c := 6.2146;

% mexp(x) ist exp(x/256); die Faktoren 256 rechnen das um.
vardef Tnorm(expr t) = mexp(-256 * c * t) enddef;

% Annahmewahrscheinlichkeit exp(-dE/T); dE in Einheiten von Tstart.
vardef Pann(expr t, dE) = mexp(-256 * dE / Tnorm(t)) enddef;

beginfig(1);
  for i = 1 upto 4:
    draw (0, i*h/4) -- (b, i*h/4)
      withcolor nxt_gray1 withpen pencircle scaled .3pt;
  endfor

  drawarrow (0,0) -- (b+12, 0) withcolor nxt_gray4 withpen pencircle scaled .5pt;
  drawarrow (0,0) -- (0, h+12) withcolor nxt_gray4 withpen pencircle scaled .5pt;

  path temp;
  temp = (0, h) for i = 1 upto 80: .. (b*i/80, h*Tnorm(i/80)) endfor;
  draw temp withcolor nxtgreen withpen pencircle scaled 1.2pt;

  numeric dE; dE := 0.12;
  path pann;
  pann = (0, h*Pann(0, dE))
    for i = 1 upto 80: .. (b*i/80, h*Pann(i/80, dE)) endfor;
  draw pann dashed evenly scaled .7 withcolor nxt_orange
    withpen pencircle scaled 1pt;

  label.rt(btex \footnotesize Temperatur $T_k=T_{\mathrm{start}}\,\alpha^k$ etex,
           (b*0.26, h*0.80)) withcolor nxtgreen;
  label.rt(btex \footnotesize Annahme einer Verschlechterung $\exp(-\Delta E/T_k)$ etex,
           (b*0.17, h*0.32)) withcolor nxt_orange;
  label.bot(btex \footnotesize Iteration $k$ etex, (b/2, -7)) withcolor nxt_gray5;
  label.lft(btex \footnotesize $T_{\mathrm{start}}$ etex, (0, h)) withcolor nxt_gray5;
  label.lft(btex \footnotesize $0$ etex, (0, 0)) withcolor nxt_gray5;
  label.bot(btex \footnotesize $0$ etex, (0, -7)) withcolor nxt_gray5;
  label.bot(btex \footnotesize $60\,000$ etex, (b, -7)) withcolor nxt_gray5;
endfig;
\end{mplibcode}

  \caption[Abkühlung und Annahmewahrscheinlichkeit]{Geometrische
    Abkühlung über den Lauf und die daraus folgende
    Annahmewahrscheinlichkeit einer Verschlechterung. Gezeichnet mit
    MetaPost.}
  \label{abb:abkuehlung}
\end{figure}

\subsection{Reproduzierbarkeit}

Simulated Annealing ist ein stochastisches Verfahren. Ein einzelner Lauf
sagt daher wenig aus, und ein Lauf ohne festen Startwert lässt sich
nicht wiederholen. Beide Punkte werden hier ernst genommen: Alle Läufe
verwenden feste Startwerte, und ausgewertet wird nicht ein Lauf, sondern
die Streuung über 20~Läufe.

\section{Werkzeuge}

Das Verfahren ist in Python~3.13 ohne weitere Abhängigkeiten
implementiert; die Auswertung nutzt NumPy \parencite{harris2020array}
und matplotlib \parencite{hunter2007matplotlib}. Als Löser dient
\ac{glpk} in Version~5.0 \parencite{glpk}. Der Quelltext des Verfahrens
findet sich in Anhang~\ref{anh:quelltext}.
